% ==============================================================================
% TSL-51: Thai Sign Language Recognition System
% ระบบจดจำภาษามือไทย TSL-51
% ==============================================================================
% บทที่ 1: บทนำ (Introduction)
% ==============================================================================

\documentclass[12pt,a4paper]{report}
\usepackage[margin=1in]{geometry}
\usepackage[utf8]{inputenc}
\usepackage[Thai]{babel}
\usepackage{graphicx}
\usepackage{hyperref}
\usepackage{amsmath}
\usepackage{amssymb}
\usepackage{algorithm}
\usepackage[noend]{algpseudocode}
\usepackage{listings}
\usepackage{xcolor}
\usepackage{longtable}
\usepackage{cite}

% Color definitions
\definecolor{codegreen}{rgb}{0,0.6,0}
\definecolor{codegray}{rgb}{0.5,0.5,0.5}
\definecolor{codepurple}{rgb}{0.58,0,0.82}
\definecolor{backcolour}{rgb}{0.95,0.95,0.92}

% Code listing style
\lstdefinestyle{mystyle}{
    backgroundcolor=\color{backcolour},
    commentstyle=\color{codegreen},
    keywordstyle=\color{blue},
    numberstyle=\tiny\color{codegray},
    stringstyle=\color{codepurple},
    basicstyle=\ttfamily\small,
    breaklines=true,
    frame=single,
    numbers=left
}
\lstset{style=mystyle}

% ==============================================================================
% DOCUMENT SETTINGS
% ==============================================================================
\title{
    \textbf{TSL-51: Thai Sign Language Recognition System} \\
    \textit{ระบบจดจำภาษามือไทย สำหรับการแปลภาษามือเป็นข้อความแบบเรียลไทม์}
}
\author{
    TSL Research Team \\
    \textit{ทีมวิจัยระบบจดจำภาษามือไทย}
}
\date{\today}

% ==============================================================================
% MAIN DOCUMENT
% ==============================================================================
\begin{document}

\maketitle
\thispagestyle{empty}

% ==============================================================================
% CHAPTER 1: บทนำ
% ==============================================================================
\chapter{บทนำ}
\pagenumbering{arabic}

\section{ความเป็นมาและความสำคัญของงานวิจัย}

ภาษามือ (Sign Language) เป็นรูปแบบการสื่อสารหลักสำหรับผู้บกพร่องทางการได้ยิน โดยในประเทศไทยมีผู้ใช้ภาษามือไทย (Thai Sign Language: TSL) จำนวนมาก 
การแปลภาษามือเป็นข้อความอัตโนมัติ (Automatic Sign Language Translation) จึงมีความสำคัญอย่างยิ่งในการช่วยให้ผู้บกพร่องทางการได้ยินสามารถสื่อสารกับผู้อื่นได้สะดวกยิ่งขึ้น 
รวมถึงการเรียนรู้ภาษามือสำหรับผู้ที่ต้องการเรียนรู้

งานวิจัยนี้นำเสนอระบบ TSL-51 (Thai Sign Language 51) สำหรับการจดจำท่าทางมือในภาษามือไทย���ำนวน 51 ท่าทาง 
ซึ่งครอบคลุมท่าทางพื้นฐานที่ใช้ในชีวิตประจำวัน ระบบที่พัฒนาขึ้นสามารถทำงานได้ทั้งในโหมดออฟไลน์ (offline) และแบบเรียลไทม์ (real-time) 
โดยใช้เทคนิค Deep Learning ร่วมกับ MediaPipe สำหรับการสกัดคุณลักษณะจากวิดีโอ

\section{วัตถุประสงค์ของงานวิจัย}

วัตถุประสงค์หลักของงานวิจัยนี้คือ:

\begin{enumerate}
    \item พัฒนาระบบจดจำภาษามือไทย TSL-51 ที่สามารถจดจำท่าทางมือได้อย่างถูกต้องและรวดเร็ว
    \item สร้างส่วนติดต่อผู้ใช้แบบ TUI (Text User Interface) สำหรับการใช้งานระบบอย่างสะดวก
    \item พัฒนาระบบแปลวิดีโอเป็นข้อความแบบเรียลไทม์โดยใช้กล้องเว็บแคม
    \item ทดสอบและเปรียบเทียบประสิทธิภาพของโมเดลต่างๆ
\end{enumerate}

\section{ขอบเขตของงานวิจัย}

งานวิจัยนี้มีขอบเขตดังนี้:

\begin{itemize}
    \item จำนวนท่าทาง: 51 ท่าทางพื้นฐานของภาษามือไทย
    \item ชุดข้อมูล: TSL-51 Dataset จาก HuggingFace
    \item โมเดลที่ใช้: GRU และ MLP (Feedforward Neural Network)
    \item การประมวลผลวิดีโอ: ใช้ MediaPipe สำหรับสกัดจุด Landmark
    \item จำนวน Feature: 162 มิติ (63 มือซ้าย + 63 มือขวา + 36 Pose)
\end{itemize}

\section{ความสำคัญของงานวิจัย}

ความสำคัญของงานวิจัยนี้มีดังนี้:

\begin{enumerate}
    \item \textbf{การช่วยเหลือผู้บกพร่องทางการได้ยิน:} ระบบที่พัฒนาขึ้นช่วยให้ผู้บกพร่องทางการได้ยินสื่อสารกับผู้อื่นได้สะดวกยิ่งขึ้น 
    โดยการแปลภาษามือเป็นข้อความแบบเรียลไทม์
    
    \item \textbf{การเรียนรู้ภาษามือ:} ระบบสามารถใช้เป็นเครื่องมือช่วยในการเรียนรู้ภาษามือไทยสำหรับผู้ที่สนใจ
    
    \item \textbf{การพัฒนาวิทยาการ:} งานวิจัยนี้เป็นพื้นฐานสำหรับการพัฒนาระบบแปลประโยคเต็ม (Sentence-level Translation) ในอนาคต
    
    \item \textbf{การเข้าถึงข้อมูล:} ชุดข้อมูล TSL-51 เปิดให้ใช้งานอย่างเป็นสาธารณะบน HuggingFace ทำให้นักวิจัยคนอื่นสามารถนำไปใช้ต่อยอดได้
\end{enumerate}

\section{การจัดโครงสร้างรายงาน}

รายงานฉบับนี้จัดโครงสร้างดังนี้:

\begin{itemize}
    \item \textbf{บทที่ 1} บทนำ - นำเสนอความเป็นมา วัตถุประสงค์ และขอบเขตของงานวิจัย
    \item \textbf{บทที่ 2} ทฤษฎีและงานวิจัยที่เกี่ยวข้อง - อธิบายทฤษฎีและงานวิจัยที่เกี่ยวข้อง
    \item \textbf{บทที่ 3} การออกแบบระบบ - อธิบายสถาปัตยกรรมและการทำงานของระบบ
    \item \textbf{บทที่ 4} การทดลองและผลลัพธ์ - นำเสนอผลการทดลองและการวิเคราะห์
    \item \textbf{บทที่ 5} บทสรุป - สรุปผลการวิจัยและข้อเสนอแนะ
\end{itemize}

% ==============================================================================
% REFERENCES
% ==============================================================================
\begin{thebibliography}{99}

bibitem{mediapipe}
MediaPipe, ``MediaPipe,'' Google, 2024. 
\url{https://mediapipe.dev}

bibitem{tsldataset}
Namonpas, ``Thai Sign Language TSL-51,'' HuggingFace, 2024.
\url{https://huggingface.co/datasets/Namonpas/thai-sign-language-tsl51}

bibitem{gru}
Cho, K., et al., ``Learning Phrase Representations using RNN Encoder-Decoder for Statistical Machine Translation,''
\textit{EMNLP 2014}.

bibitem{pytorch}
Paszke, A., et al., ``PyTorch: An Imperative Style, High-Performance Deep Learning Library,''
\textit{NeurIPS 2019}.

bibitem{mlkit}
Keras, ``Module,'' TensorFlow, 2024.
\url{https://keras.io/api/losses/}

bibitem{cvpr}
Zhong, Y., et al., ``Camera Gesture Recognition with Deep Learning,'' 
\textit{arXiv preprint arXiv:1705.10223}, 2017.

bibitem{tslavatar}
Wiwat, J., et al., ``Thai Sign Language Recognition using Deep Learning,'' 
\textit{ICEE 2020}.

\end{thebibliography}

\end{document}